\def\stat{krivenko}





\def\tit{СТАТИСТИЧЕСКИЙ КРИТЕРИЙ СТАБИЛЬНОСТИ СИСТЕМЫ МАССОВОГО ОБСЛУЖИВАНИЯ, 


ОСНОВАННЫЙ НА~ПОСЛЕДОВАТЕЛЬНОСТИ ВРЕМЕН ПРЕБЫВАНИЯ}





\def\titkol{Статистический критерий стабильности системы массового обслуживания}


%,  основанный на~последовательности времен пребывания}





\def\aut{М.\,П.~Кривенко$^1$}





\def\autkol{М.\,П.~Кривенко}





\titel{\tit}{\aut}{\autkol}{\titkol}





\index{Кривенко М.\,П.}


\index{Krivenko M.\,P.}





%{\renewcommand{\thefootnote}{\fnsymbol{footnote}}


%\footnotetext[1]


%{Работа выполнялась с~использованием инфраструктуры Центра коллективного пользования 


%<<Высокопроизводительные вычисления и~большие данные>> (ЦКП <<Информатика>>) ФИЦ ИУ РАН 


%(г.~Москва).}} 





\renewcommand{\thefootnote}{\arabic{footnote}}


\footnotetext[1]{Федеральный исследовательский центр <<Информатика и~управление>> Российской академии наук, 


\mbox{mkrivenko@ipiran.ru}}








 \label{st\stat}





 \thispagestyle{headings}


 


\vspace*{-12pt}

















      \Abst{Рассматривается задача статистической проверки ста\-биль\-ности сис\-тем массового 


обслуживания (СМО) на основе характеристик по\-сле\-до\-ва\-тель\-ности значений времени 


пребывания заявок в~сис\-те\-ме. Для ее решения предлагается использовать критерии единичного 


корня. Устанавливается связь между количественными и~временными характеристиками 


сис\-те\-мы в~стиле формулы Литтла для нестационарного случая, что позволяет сопоставить  


rate-ста\-биль\-ность со ста\-ци\-о\-нар\-ностью по\-сле\-до\-ва\-тель\-ности времен пребывания. В~качестве 


критериев ста\-ци\-о\-нар\-ности рассматриваются три базовые модели и~соответствующие тесты 


Ди\-ки--Фул\-ле\-ра. Об\-суж\-да\-ют\-ся ограничения применяемых статистических таб\-лиц. Проведенное 


имитационное исследование поз\-во\-ля\-ет сделать выводы: все три теста решают по\-став\-лен\-ную 


задачу анализа ста\-биль\-ности с~явной потерей качества при\-ни\-ма\-емых решений в~случае модели 


с~детерминированным линейным трендом, в~об\-ласти ста\-биль\-ности сист\-ем критерии ведут себя 


хуже, чем в~об\-ласти не\-ста\-биль\-ности; явно проявляются ожи\-да\-емые тенденции применения 


тес\-тов; некоторые отклонения от ожи\-да\-емых результатов говорят о~не\-об\-хо\-ди\-мости расширять 


спектр при\-ме\-ня\-емых моделей и~снимать ограничения ис\-поль\-зу\-емых статистических таб\-лиц.}


      


      \KW{система массового обслуживания; анализ временных рядов; тесты единичного 


корня; тесты Ди\-ки--Фул\-ле\-ра; имитационное моделирование}


      


\DOI{10.14357\!/\!08696527240204}{RVWWXQ}  





\vskip 10pt plus 9pt minus 6pt





 \thispagestyle{headings}


 


 \vspace*{-12pt}





\section{Введение}





\vspace*{-4pt}








     Функционирование СМО предполагает, что через ее вход и~выход 


соответственно прибывают и~убывают заявки. Для каждого неотрицательного 


времени~$t$ пусть $A(t)$ обозначает чис\-ло прибывших заявок (входной процесс), 


$D(t)$~--- покинувших сис\-те\-му (выходной процесс). Тогда общее чис\-ло~$N(t)$~заявок в~сис\-те\-ме в~момент~$t$ равно


     $$


     N(t)=N(0)+A(t)-D(t).


     $$


Данная общая модель вход\-но\-го--вы\-ход\-но\-го процесса~$N(t)$ 


подразумевает сле\-ду\-ющие свойства: $A(t)$ и~$D(t)$ суть неубывающие 


непрерывные справа процессы, $N(t)\hm\geq 0$ для всех $t\hm\geq 0$. 


     


     Для времени прибытия и~убытия $k$-й заявки, $k\hm=1,2,\ldots$, введем 


самостоятельные обозначения~$A_k$ и~$D_k$. Тогда 


     $$ 


     A(t)=\# \left\{ k: A_k\leq t\right\};\quad D(t)=\# \left\{ k: D_k\leq t\right\}.


     $$


     %


     Также определим:


         $N(t)= \# \{ k: A_k\leq t< D_k\}\hm= A(t)\hm-D(t)$~---  чис\-ло заявок  


в~сис\-те\-ме;


          $V_k=D_k-A_k$~--- время пребывания $k$-й заявки в~сис\-теме.


     


     Пусть $1\{E\}$~--- индикатор события~$E$. Тогда


     \begin{align*}


     N(t) &=\displaystyle \sum\limits^\infty_{k=1} 1 \left\{ A_k\leq t < D_k\right\};\\


     V_k&= \int\limits_0^\infty 1\left\{ A_k\leq t<D_k\right\}dt\,.


     \end{align*}





\section{Связь определений стабильности}





     Ранее стабильность системы со ссылками на основополагающие работы 


была определена в~[1] через поведение вход\-но\-го--вы\-ход\-но\-го 


процесса~$N(t)$ или, что эквивалентно, на основе входного и~выходного 


процессов~$A(t)$ и~$D(t)$. Если требуется сделать это с~по\-мощью временн\'{о}го 


ряда $\{ V_k,\ k\hm=1,2,\ldots\}$, то необходимо установить связь между 


количественными и~временн\'{ы}ми характеристиками СМО, например в~стиле 


формулировок закона Литтла, но при отказе от условия ста\-ци\-о\-нар\-ности. 





Заменим 


равенство Литтла на ограничения снизу и~сверху для $\int_0^t L(s)\,ds$ на основе 


представления





\vspace*{2pt}





\noindent


     $$


     \int\limits_0^t N(s)\,ds =\int\limits_0^t \sum\limits_{k=1}^\infty 1\{ A_k\leq 


s<D_k\}\,ds\,.


     $$


     %


     Сначала после отбрасывания час\-ти сла\-га\-емых имеем:


     


     \vspace*{-6pt}


     


     \noindent


     \begin{multline*}


     \int\limits_0^t \sum\limits^\infty_{k=1} 1\left\{ A_k\leq s<D_k\right\}ds \geq 


     \int\limits_0^t \sum\limits_{\{k:D_k<t\}} 1\left\{ A_k\leq s<D_k\right\}ds={}\\


     {}= 


     \!\!\!\sum\limits_{\{ k:D_k<t\}} \int\limits_0^t 1\left\{ A_k\leq s<D_k\right\}ds= 


\!\!\!\sum\limits_{\{k: D_k<t\}} \int\limits_0^\infty 1\left\{ A_k\leq s< 


D_k\right\}ds=\!\!\!\sum\limits_{\{k: D_k<t\}}\!\!\!\! V_k.


     \end{multline*}


     %


     Далее за счет расширения об\-ласти интегрирования получаем:


     


     \noindent


     \begin{multline*}


     \int\limits_0^t \sum\limits^\infty_{k=1} 1\left\{ A_k\leq s<D_k\right\}ds= {}\\


     {}=


     \int\limits_0^t\left[ \sum\limits_{\{k: A_k<t\}} 1\left\{ A_k\leq s<D_k\right\} + 


\sum\limits_{\{k: A_k\geq t\}} 1\left\{ A_k\leq s<D_k\right\} \right] ds={}\\


     {}= \!\!\int\limits_0^t \sum\limits_{\{k: A_k<t\}} 1\left\{ A_k\leq s<D_k\right\} ds 


\leq


     \!\!\int\limits_0^\infty \sum\limits_{\{k: A_k<t\}} 1\left\{ A_k\leq s<D_k\right\} ds 


=\!\!\!\sum\limits_{\{k: A_k<t\}}\!\!\! V_k.\hspace*{-2.42pt}


     \end{multline*}


Таким образом приходим к~базовым соотношениям:


$$


\sum\limits_{\{k: D_k<t\}} V_k\leq \int\limits_0^t N(s)\,ds\leq \sum\limits_{\{k: 


A_k<t\}} V_k.


$$


После деления всех частей полученных неравенств на~$t$ и~использования 


пред\-став\-ле\-ний вида


$$


\fr{\sum\nolimits_{\{k: \ldots\}} V_k}{t}=\fr{\#\{k:\ldots\}}{t}\, 


\fr{\sum\nolimits_{\{k:\ldots\}} V_k}{\#\{k:\ldots\}}


$$


получаем элементы формулы Литтла и~важные для статистического контроля 


стабильности выводы: для стабильного случая время пребывания заявки 


в~сис\-те\-ме в~среднем постоянно, а~для нестабильного~--- растет по крайней мере линейно.





     Обратим внимание, что левая и~правая час\-ти полученных базовых 


соотношений становятся равными в~периоды про\-стоя, т.\,е.\ для таких моментов 


времени~$t$, когда $N(t)\hm=0$. Для стабильных сис\-тем наличие таковых 


возможно, но, строго говоря, ниоткуда не следует. Поэтому при использовании 


равенств, например для обоснования формулы Литтла для стабильных сис\-тем, 


в~[2] пришлось дополнительно вводить условие существования периодов про\-стоя.


     


     С~точки зрения статистического контроля функционирования СМО важ\-нее 


другое: для нестабильных сис\-тем~$N(t)$ линейно растет при 


уве\-ли\-чи\-ва\-ющем\-ся~$t$, поэтому время пребывания в~сис\-те\-ме в~среднем также 


растет по крайней мере линейно, т.\,е.\ не\-ста\-биль\-ность на основе входного 


и~выходного потоков по\-рож\-да\-ет изменения характеристик по\-сле\-до\-ва\-тель\-ности 


времен пребывания в~сис\-теме. 





\section{Тесты единичного корня анализа стационарности}





     Тематика <<единичные корней>> относится к~использованию общей 


линейной модели (ОЛМ) при анализе временн\'{ы}х рядов в~первую очередь  


в~об\-ласти эко\-но\-мет\-ри\-ки, ее применение также отмечается~[3] в~отдельных 


отраслях геонауки, в~физиологии и~медицинской диагностике, при 


социологических исследованиях, в~ходе анализа веб-тра\-фи\-ка. Вместе с~тем 


пуб\-ли\-ка\-ции, посвященные тестированию единичного корня применительно 


к~СМО, автору \mbox{статьи} найти не удалось и~в~первую очередь из-за 


ма\-ло\-чис\-лен\-ности работ по статистическому контролю функционирования таких 


сис\-тем в~целом. 


     


     Как оказывается, анализируемые последовательности времен пребывания 


обладают специфическими свойствами~[1]: высокой сте\-пенью временн\'{о}й 


свя\-зан\-ности и~наличием или отсутствием ста\-ци\-о\-нар\-ности. Эти свойства 


интересны как особенности временн\'{ы}х рядов вообще и~актуальны с~точки зрения 


установления ста\-биль\-ности. Последнее приобретает первостепенное значение при 


реализации методологии Бок\-са--Джен\-кин\-са: если, например, в~об\-ласти 


эко\-но\-мет\-ри\-ки наличие тренда обычно становится ме\-ша\-ющим фактором при 


решении задач прогнозирования, то при установлении ста\-биль\-ности СМО 


выявление роста сред\-не\-го времени пребывания становится ключевым. При этом 


выбор адекватной модели должен не только приводить к~непротиворечивым 


оценкам ее па\-ра\-мет\-ров и~сводить к~минимуму ошибочные решения в~первую 


очередь при контроле ста\-биль\-ности, но и~по воз\-мож\-ности формировать основу 


решения иных задач анализа данных, как, например, прогноза. 


     


     Тестирование единичного корня предполагает, что ис\-сле\-ду\-емый временной 


ряд $\{y_k,\ k\hm=1,2,\ldots\}$ может быть пред\-став\-лен как


     $$


     y_k=d_k+s_k+\varepsilon_k\,,


     $$


где $d_k$~--- детерминированный компонент (постоянное смещение, 


функционально описанный тренд, сезонные изменения и~т.\,п.); $s_k$~--- 


стохастический компонент (процессы авторегрессии~--- скользящего среднего 


общего и~част\-но\-го вида); $\varepsilon_k$~--- стационарный процесс ошибок. 


Задача состоит в~том, чтобы определить, является~$s_k$ нестационарным или 


стационарным процессом. 





     Стационарность далее понимается в~широком смысле, когда математическое 


ожидание случайного процесса постоянно, а~корреляционная функция зависит 


только от раз\-ности аргументов. Интегрированный  


(раз\-ност\-но-ста\-ци\-о\-нар\-ный) процесс $I(n)$ порядка~$n$~--- это тот, для 


которого надо перейти к~разностям~$n$~раз, чтобы он стал стационарным. Для 


обычным образом определенной модели авторегрессии~--- сколь\-зя\-ще\-го сред\-не\-го 


и~характеристического полинома~$\varphi(z)$ ее авторегрессионной части 


$\mathrm{AR}\,(p)$~--- ста\-ци\-о\-нар\-ность моделируемого процесса определяется свойствами 


корней $\varphi(z)$: они все долж\-ны быть вне единичного круга~[4].


     


     Процесс единичного корня определяется через  $\mathrm{AR}\,(1)$ при наличии единичного по 


модулю корня~$z_1$, при этом внимание сосредоточивается на $z_1\hm=1$. 


Когда $z_1\hm=-1$, процесс будет осциллировать (знак меняется на 


противоположный при переходе от одного элемента временн\'{о}го ряда к~другому), 


что не соответствует ни представлению о принципах функционирования СМО, ни 


име\-ющим\-ся результатам наблюдений времен пребывания заданий. Для 


единичного корня временной ряд принимает вид $y_k\hm=  


y_{k-1}\hm+\varepsilon_k$ и~становится интегрированным процессом первого 


порядка. 


     


     Критерии единичного корня нацелены на нулевую гипотезу, состоящую 


в~том, что наблюдаемый временной ряд~--- интегрированный процесс 1-го 


порядка $I(1)$. Базовая модель временн\'{о}го ряда пред\-став\-ля\-ет собой $\mathrm{AR}\,(1)$ 


с~воз\-мож\-ностью пополнения ее постоянной составляющей и~линейным 


трендом и~приводит к~тестам Ди\-ки--Фул\-ле\-ра (Dickey--Fuller, DF). 


Расширенный вариант $\mathrm{DF}$-тес\-тов использует модель авторегрессии~--- 


скользящего среднего с~компонентами разностей, после <<удаления>> которых во 


вспомогательной регрессии он может использовать критические значения  


$\mathrm{DF}$-тес\-та. Кон\-ку\-ри\-ру\-ющая гипотеза ориентирована на предположение 


о~ста\-ци\-о\-нар\-ности.


     


     Перечисленные критерии не исчерпывают все существующие варианты 


тес\-тов единичного корня~[5], их сравнительный анализ на основе трех основных 


аспектов~--- мощ\-ности тес\-та, его обосно\-ван\-ности и~прос\-то\-ты использования был 


проведен в~[6]. Подобное имитационное комплексное исследование относится 


к~редким и~достаточно узко направленным (в~частности, рассмотрение только 


модели~$\mathrm{AR}\,(1)$ и~только с~положительным значением па\-ра\-мет\-ра), но оно может 


стать отправной точ\-кой при выборе критерия стабильности СМО на осно\-ве 


единичного корня в~пользу $\mathrm{DF}$-тес\-тов. 


     


     Базовый вариант $\mathrm{DF}$-тестов основывается на представлении


     \begin{equation}


     y_k=\rho y_{k-1}+\varepsilon_k\,,\enskip k=1,2,\ldots, \enskip y_0=0\,.


     \label{e1-kri}


     \end{equation}


Нулевая гипотеза $H_0: \rho =1$, т.\,е.\ наблюдаемый процесс~--- $I(1)$ 


и~поэтому нестационарный. Конкурирующая гипотеза $H_1: \vert\rho\vert \hm<1$, 


т.\,е.\ речь идет о~$I(0)$ и~стационарности.





     Модель~(\ref{e1-kri}) может быть расширена путем до\-бав\-ле\-ния константы 


     \begin{equation}


     y_k= \mu +\rho y_{k-1} +\varepsilon_k\,.


     \label{e2-kri}


     \end{equation}


или линейного тренда


\begin{equation}


y_k= \mu+\beta k+ \rho y_{k-1} +\varepsilon_k\,.


\label{e3-kri}


\end{equation}


     


     Тестовая статистика имеет вид


     $(\hat{\rho}\hm-1)/\mathrm{SE}\,(\hat{\rho})$,


где $\hat{\rho}$~--- оценка наименьших квадратов коэффициента~$\rho$; 


$SE(\hat{\rho})$~--- ее стандартная ошибка для регрессионной модели; $p$-зна\-че\-ния


 табулированы методом статистических испытаний для стандартных 


процентилей.


     


     Принимая во внимание поведение $\{V_k\}$ для стабильной и~нестабильной 


сис\-тем, рас\-смот\-рим по\-дроб\-нее модели~(2) и~(\ref{e3-kri}).


     


     Модель (\ref{e2-kri}) при предположениях $\mu\not= 0$ и~$\rho\hm=1$ 


можно переписать как


     \begin{equation}


     y_k= y_0 +\mu k+E_k\,, 


     \label{e4-kri}


     \end{equation}


где $E_k=\sum\nolimits^k_{i=1} \varepsilon_i$.


Тогда для нестабильной системы асимп\-то\-ти\-че\-ски получаем процесс 


с~до\-ми\-ни\-ру\-ющим линейным рос\-том, так как  по ве\-ро\-ят\-ности $E_k\hm= O(k^{1/2})$. 


Отсюда видно, что по крайней мере общий характер поведения сред\-не\-го $\{V_k\}$ 


получает в~этой модели свое отражение. В~случае стабильной сис\-те\-мы, для 


которой $\rho\hm\in [0,1)$, по аналогии с~(\ref{e4-kri}) получаем


$$


y_k= \rho^k y_0 +\mu \sum\limits^k_{i=0} \rho^i +\sum\limits^k_{i=0} \rho^i 


\varepsilon_{k-i}.


$$


Это приводит к~постоянству среднего для $\{V_k\}$, которое за счет 


выбора~$y_0$ и~$\mu$ может стать нулевым, т.\,е.\ опять же получаем нуж\-ное 


поведение исследуемого временн\'{о}го ряда. 





     Естественной альтернативной моделью для $\{V_k\}$ может стать 


комбинация линейного тренда и~стационарного временн\'{о}го ряда:


     \begin{equation}


     y_k= \beta_0 +\beta_1  k +x_k\,,


     \label{e5-kri}


     \end{equation}


где $x_k= \rho x_{k-1} +\varepsilon_k$. После подстановки авторегрессии 


в~(\ref{e5-kri}) получаем пред\-став\-ле\-ние вида~(\ref{e3-kri}), а~именно:


$$


y_k= \gamma_0 + \gamma_1 k +\rho y_{k-1} +\varepsilon_k\,,


$$


     где 


     $$


     \gamma_0= \beta_0(1-\rho)+\rho \beta_1\,;\quad \gamma_1= \beta_1(1-\rho)\,.


     $$


Модель~(\ref{e5-kri}) включает детерминированный линейный временной 


тренд, не связанный со значением~$\rho$. Если $\rho\hm=1$, то $\gamma_0\hm= 


\beta_1$, $\gamma_1\hm=0$ и~$y_k\hm= \beta_1\hm+ y_{k-1} \hm+ \varepsilon_k$, 


т.\,е.\ тренд~(\ref{e5-kri}) с~наложенной на него авторегрессией фактически дает 


$\mathrm{AR}\,(1)$ с~константой. 





\section{Таблицы тестов}





     Набор таблиц встречается в~ряде изданий, но не всегда в~одинаковом виде 


     и~с~достаточными обоснованиями. Подробное описание по\-стро\-ения таб\-лиц 


     с~оценкой по\-греш\-ности приведено в~[7], вместе с~отдельными уточ\-не\-ни\-ями их 


мож\-но найти в~приложении~A~[8]. Более важным, чем незначительные 


отклонения в~табличных данных, оказывается то, что обычно не приводятся 


условия, при которых они получены. 


     


     Так, для модели~(2) предельное распределение статистики табулировано 


в~предположении, что $\mu\hm=0$, а~для~(3)~--- $\beta\hm=0$ (значение~$\mu$ 


в~модели~(\ref{e3-kri}) на распределение не влияет). Если $\mu\not= 0$ 


для~(\ref{e2-kri}) или $\beta\not= 0$ для~(\ref{e3-kri}), то предельные 


распределения статистик критерия становятся нормальными с~со\-от\-вет\-ст\-ву\-ющи\-ми 


моментами~[7], что объясняет следующее: если приняты описанные модели~(2) 


и~(3), а для поверки гипотезы $\rho\hm=1$ используются со\-от\-вет\-ст\-ву\-ющие 


статистики, то гипотеза будет принята с~ве\-ро\-ят\-ностью, большей, чем та, которая 


соответствует нулевым значениям~$\mu$ и~$\beta$. Эту специфику использования 


таб\-лиц необходимо принимать во внимание при верификации со\-от\-вет\-ст\-ву\-юще\-го 


программного обеспечения, а~также в~ходе сравнительного анализа мощ\-ности 


критериев зна\-чи\-мости.


     


     Следует обратить внимание еще на одну деталь применения критериев на 


практике. В~приведенных рас\-суж\-де\-ни\-ях начальное значение фиксировано 


и~обычно $y_0\hm=0$. Но если предельные распределения статистик критериев 


не зависят от~$y_0$, то для малых объемов данных это не так. Актуальным это 


может стать при разделении наблюдаемых процессов на части с~по\-сле\-ду\-ющим 


объединением результатов их статистического анализа (подобно тому, как 


вводятся в~рас\-смот\-ре\-ние фазы и~фрагменты в~[1]). Для того чтобы снизить 


последствия указанной за\-ви\-си\-мости, до\-ста\-точ\-но прос\-то все элементы 


временн\'{о}го ряда из фрагмента сдвинуть по значению на величину начального 


(первого) элемента, в~результате чего оно станет нулевым. Не надо забывать, что 


цель со\-сто\-ит в~контроле стационарности временн\'{о}го ряда $\{V_k\}$, 


а~одинаковый сдвиг всех наблюденных значений в~этом смыс\-ле не влияет на 


результаты. 


     


     И, наконец, предельные распределения статистик справедливы для ошибок, 


которые пред\-став\-ля\-ют собой независимые и~одинаково распределенные 


(необязательно нормально) случайные величины с~нулевым средним и~некоторой 


дисперсией~[9].


     


\section{Мощность тестов единичного корня}





     В~условиях многовариантности тестов единичного корня и~отсутствия 


практики их применения для статистического контроля функционирования СМО 


становится актуальным сравнение пред\-ла\-га\-емых подходов и~решений. При этом 


надо помнить, что при решении задачи выявления не\-ста\-биль\-ности попутно 


реализуется фундаментальный этап по\-стро\-ения модели про\-те\-ка\-ющих процессов. 


     


     Сосредоточив внимание на трех описанных моделях для $\{V_k\}$, 


продолжим рассмотрение задачи контроля стабильной работы двухпроцессорной 


сис\-те\-мы обработки заданий со случайным выбором числа требуемых 


процессоров~[1]. Функционирование подобной сис\-те\-мы $M/M/2$ определяется 


параметрами: $\lambda$~--- ин\-тен\-сив\-ность по\-ступ\-ле\-ния заданий на обработку; 


$p_1$~--- ве\-ро\-ят\-ность того, что для выполнения задания требуется один 


процессор; $\mu$~--- сред\-нее время обработки задания  


про\-цес\-со\-ром\!/\!про\-цес\-со\-ра\-ми. Условие ста\-биль\-ности 


функционирования этой сис\-те\-мы принимает вид $\lambda/\mu \hm< 2/(2\hm- p_1^2)$, 


с~по\-мощью которого будет осуществляться выбор мо\-де\-ли\-ру\-емых  


сис\-тем: па\-ра\-мет\-ры $p_1$ и~$\mu$ принимаются базовыми со значениями 


$p_1\hm= 0{,}5$ и~$\mu\hm=1$, а~$\lambda$ выбирается как $C_\lambda \cdot 2\mu/(2-p_1^2)$, 


где $C_\lambda$~--- индикатор ста\-биль\-ности работы сис\-те\-мы (при 


$C_\lambda\hm<1$ сис\-те\-ма стабильна). Имитационное исследование проводилось 


для фрагментов ряда $\{V_k\}$, на\-чи\-на\-ющих\-ся с~$k\hm\approx 5000$ (чтобы 


исключить влияние начального этапа функционирования СМО) и~име\-ющих 


размер из набора $n\hm= 100$, 200, 400, 800 (чтобы проследить влияние 


увеличения объема анализируемых данных), и~включало~1000~экспериментов. 


Уровень зна\-чи\-мости~$\alpha$ для $\mathrm{DF}$-тес\-тов принимался рав\-ным~5\% 


или~10\% (чтобы выяснить степень важ\-ности использования определенного 


значения~$\alpha$ при принятии решений). Иллюстрацией характера полученных 


результатов может служить рисунок, на котором для $n\hm=800$ и~$\alpha\hm= 


10\%$ приведены графики частоты правильных решений $\mathrm{DF}$-тес\-тов для 


различных вариантов сис\-темы. 


     


     


\begin{figure} %fig1


      \vspace*{1pt}


  \begin{center}


 \mbox{%


 \epsfxsize=62mm  


\epsfbox{kri-1.eps}


 }


\end{center}


%\vspace*{-12pt}





\noindent


{\small Зависимость частоты правильных решений $\mathrm{DF}$-тестов для СМО с~различными 


значениями индикатора ста\-биль\-ности~$C_\lambda$: \textit{1}~--- модель~(1); \textit{2}~--- 


модель~~(2); \textit{3}~--- модель~(3)}


\end{figure}





     Из результатов моделирования можно сделать общие выводы:


     \begin{itemize}


\item все три критерия решают по\-став\-лен\-ную задачу анализа ста\-биль\-ности с~явной 


потерей качества при\-ни\-ма\-емых решений в~случае модели~(3);


\item в~области ста\-биль\-ности ($C_\lambda\hm<1$) критерии ведут себя хуже, чем 


в~об\-ласти нестабильности, рост ошибок наблюдается при при\-бли\-же\-нии к~границе 


$C_\lambda\hm=1$;


\item явно проявляются ожидаемые тенденции применения тестов, а~именно: рост 


пра\-виль\-ности при\-ни\-ма\-емых решений с~увеличением объема ана\-ли\-зи\-ру\-емых 


данных; рост частоты принятия гипотезы о~не\-ста\-биль\-ности с~уменьшением уровня 


значимости (что приводит к~росту ошибок при анализе стабильных сис\-тем);


\item некоторые отклонения от ожи\-да\-емых результатов (например, конкретные 


значения~$\alpha$ иногда не отвечают результатам моделирования) говорят 


о~неполной адекват\-ности при\-ни\-ма\-емых моделей данных и~ограниченных 


возможностях таб\-лиц DF-тес\-тов.


\end{itemize}





\section{Заключение}





     Установленная связь между количественными и~временн\'{ы}ми 


характеристиками СМО для нестационарного случая поз\-во\-ли\-ла со\-по\-ста\-вить rate-ста\-биль\-ность 


со ста\-ци\-о\-нар\-ностью ряда времен пребывания заявок в~сис\-те\-ме. Это 


позволяет расширить набор статистических процедур контроля ста\-биль\-ности за 


счет критериев стационарности типа тестов Ди\-ки--Фул\-ле\-ра. Результаты 


экспериментов показали, что они остаются надежным вариантом для тестирования 


единичного корня, особенно в~случае временн\'{ы}х рядов с~большим чис\-лом 


наблюдений. Также они остаются одними из наиболее час\-то ис\-поль\-зу\-емых тес\-тов, 


поскольку их решающее преимущество заключается в~прос\-то\-те конструкции и~осу\-щест\-ви\-мости. 


     


     Таким образом, проведенное исследование создает предпосылки для 


существенного расширения в~будущем как спект\-ра при\-вле\-ка\-емых моделей, так 


и~собственно методов анализа данных. 


     


{\small\frenchspacing


{%\baselineskip=10.8pt


\addcontentsline{toc}{section}{Литература}


\begin{thebibliography}{9}


\bibitem{1-kri}


\Au{Кривенко М.\,П.} Статистический критерий ста\-биль\-ности сис\-те\-мы массового 


обслуживания, основанный на входном и~выходном потоках~/\!\!/


 Информатика и~её применения, 2024. Т.~18. Вып.~1. С.~54--60. doi: 10.14357\!/\!19922264240108. EDN: 


JNJJMU.


\bibitem{2-kri}


\Au{El-Taha~M., Stidham~S., Jr.} Sample-path analysis of queueing systems.~--- New 


York, NY, USA: Springer, 1999. 302~p. doi:  


10.1007\!/\!978-1-4615-5721-0.


\bibitem{3-kri}


\Au{Herranz~E.} Unit root tests~/\!\!/ Wiley Interdisciplinary Reviews 


Computational Statistics, 2017. Vol.~9. Iss.~3. 


Art.~e1396.  20~p. doi: 10.1002\!/\!wics.1396.


\bibitem{4-kri}


\Au{Chan N.\,H.} Time series: Applications to finance.~--- New York, NY, USA: 


Wiley, 2002. 225~p.


\bibitem{5-kri}


\Au{Kocenda E., Cern\'{y}~A.} Elements of time series econometrics: An applied 


approach.~--- Prague, Czech Republic: Charles University, Karolinum Press, 2007. 


228~p.


\bibitem{6-kri}


\Au{Arltova M., Fedorova~D.} Selection of unit root test on the basis of length of the 


time series and value of $\mathrm{AR}\,(1)$ parameter~/\!\!/ Statistika Statistics Economy~J., 


2016. Vol.~96. No.\,3. P.~47--64.


\bibitem{7-kri}


\Au{Dickey D.\,A.} Estimation and hypothesis testing in nonstationary time series: PhD 


Diss.~--- Ames, IA, USA: Iowa State University Digital Repository, 1976. 


128~p. doi: 


10.31274\!/\!rtd-180813-2848. 


%https://dr.lib.iastate.edu/entities/publication/0cdcd98a-18aa-4b11-9b7b-44a5acef6560.


\bibitem{8-kri}


\Au{Fuller W.\,A.} Introduction to statistical time series.~--- 2nd ed.~--- Hoboken, NJ, 


USA: Wiley, 1996. 728~p. doi: 10.1002\!/\!9780470316917.


\bibitem{9-kri}


\Au{Hasza D.\,P.} Estimation in nonstationary time series: PhD Diss.~--- Ames, 


IA, USA: Iowa State University Digital Repository, 1977. 130~p.


doi: 10.31274\!/\!rtd-180813-5232.


% {\sf  https://dr.lib.iastate.edu/entities/publication/6c932b69-4f78-4647-9ac4-068b3acca90c}.


     \end{thebibliography}


} }











\vspace*{-6pt}





\hfill{\small\textit{Поступила в~редакцию 15.03.24}}





%\vspace*{8pt}





%\hrule





%\vspace*{2pt}








%\hrule





%\vspace*{8pt}





\newpage





\vspace*{-28pt}





\def\tit{STATISTICAL CRITERION FOR~QUEUING SYSTEM STABILITY 


BASED~ON~TIME SPENT SERIES}











\def\titkol{Statistical criterion for~queuing system stability 


based~on~time spent series}





\def\aut{M.\,P.~Krivenko}





\def\autkol{M.\,P.~Krivenko}





\titel{\tit}{\aut}{\autkol}{\titkol}





\vspace*{-8pt}





\noindent


Federal Research Center ``Computer Science and Control'' of the Russian Academy


of Sciences,   44-2~Vavilov Str., Moscow 119133, Russian Federation








\def\leftfootline{\small{\textbf{\thepage}


\hfill Sistemy i Sredstva Informatiki~---


Systems and Means of Informatics 2024 vol~34 no\,2}


}%


 \def\rightfootline{\small{Sistemy i Sredstva Informatiki~---


 Systems and Means of Informatics 2024 vol~34 no\,2


\hfill \textbf{\thepage}}}





\vspace*{3pt}











\Abste{The article discusses the problem of statistical test of queuing system stability based on 


the characteristics of the time spent series  $\{ V_k\}$. 


To solve it, it is proposed to use the tests for unit roots. A~relationship is established between the quantitative and temporal 


characteristics of the system in the style of the Little's formula for the nonstationary case which makes it possible to 


link the rate-stability with the stationarity  $\{V_k\}$. Three basic models and the corresponding Dickey--Fuller 


tests are considered as the stationarity test. The limitations of the used statistical tables are discussed. The 


conducted simulation study allows one to draw the following conclusions: all three criteria solve the problem of stability 


analysis with a clear loss of the quality of decisions made in the case of a~model with a~deterministic linear trend, 


in the field of system stability, the test behaves worse than in the field of instability; the expected trends 


in the use of tests are clearly manifested; and some deviations from 


the expected results indicate the need to expand the range of the models used and to remove the limitations of the statistical tables used.}





\KWE{queueing system; time series analysis; unit root tests; Dickey--Fuller tests; 


simulation}





\DOI{10.14357\!/\!08696527240204}{RVWWXQ}  





%\vspace*{-10pt}








%\Ack


%\vspace*{-3pt}


%\noindent








%\vspace*{-3pt}





\renewcommand{\bibname}{\protect\rmfamily References}


%\renewcommand{\bibname}{\large\protect\rm References}





{\small\frenchspacing


{ %\baselineskip=12pt


\addcontentsline{toc}{section}{References}


\begin{thebibliography}{9}


\bibitem{1-kri-1}


\Aue{Krivenko, M.\,P.} 2024. Sta\-ti\-sti\-che\-skiy kri\-te\-riy sta\-bil'\-nosti sis\-te\-my mas\-so\-vo\-go 


ob\-slu\-zhi\-va\-niya, osno\-van\-nyy na vkhod\-nom i~vy\-khod\-nom po\-to\-kakh [Statistical criterion 


for queuing system stability based on input and output flows]. \textit{Informatika i~ee 


Primeneniya~--- Inform. Appl.} 18(1):54--60. doi: 10.14357\!/\!19922264240108. EDN: 


JNJJMU.


\bibitem{2-kri-1}


\Aue{El-Taha, M., and S.~Stidham, Jr.} 1999. \textit{Sample-path analysis of queueing 


systems}. New York, NY: Springer. 302~p. doi: 10.1007\!/\!978-1-4615-5721-0.


\bibitem{3-kri-1}


\Aue{Herranz, E.} 2017. Unit root tests. \textit{Wiley Interdisciplinary Reviews 


Computational Statistics} 9(3):e1396. 20~p. doi: 10.1002\!/\!wics.1396.


\bibitem{4-kri-1}


\Aue{Chan, N.\,H.} 2002. \textit{Time series: Applications to finance}. New York, NY: 


Wiley. 225~p.


\bibitem{5-kri-1}


\Aue{Kocenda, E., and A.~Cern\'{y}}. 2007. \textit{Elements of time series 


econometrics: An applied approach}. Prague, Czech Republic: Charles University, 


Karolinum Press. 228~p.


\bibitem{6-kri-1}


\Aue{Arltova, M., and D.~Fedorova.} 2016. Selection of unit root test on the basis of 


length of the time series and value of $AR(1)$ parameter. \textit{Statistika Statistics 


Economy~J.} 96(3):47--64.


\bibitem{7-kri-1}


\Aue{Dickey, D.\,A.} 1976. Estimation and hypothesis testing in nonstationary time 


series.  Ames, IA: Iowa State University Digital Repository.  PhD Diss. 128~p. doi: 


10.31274\!/ rtd-180813-2848. 


\bibitem{8-kri-1}


\Aue{Fuller, W.\,A.} 1996. \textit{Introduction to statistical time series}. 2nd ed. 


Hoboken, NJ: Wiley. 728 p. doi: 10.1002\!/\!9780470316917.


\bibitem{9-kri-1}


\Aue{Hasza, D.\,P.} 1977. Estimation in nonstationary time series.  Ames, IA: 


Iowa State University Digital Repository.  PhD Diss. 130~p. doi: 10.31274\!/\!rtd-180813-5232.





\end{thebibliography}


} }





\vspace*{-4pt}





\hfill{\small\textit{Received March 15, 2024}} 





\vspace*{-10pt} 





\Contrl





\vspace*{-3pt}





\noindent


\textbf{Krivenko Michail P.} (b.\ 1946)~--- Doctor of Science in technology, professor, 


leading scientist, Federal Research Center ``Computer Science and Control'' of the 


Russian Academy of Sciences, 44-2~Vavilov Str., Moscow 119333, Russian Federation; 


\mbox{mkrivenko@ipiran.ru}








\label{end\stat}





\renewcommand{\bibname}{\protect\rm Литература} 


